\section{Expériences de recherche}

        \begin{twocolentry}{
            août 2025 - fév. 2026
        }
            \textbf{Quantifier l'invisible : Les femmes médecins dans les Guides Rosenwald}

            [\href{https://medif.hypotheses.org/2006}{Blog}] [\href{https://renyi.ch/reports/rosenwald.pdf}{Rapport}] [\href{https://arxiv.org/abs/2605.25781}{arXiv}] [\href{https://github.com/nmrenyi/double-triangle-annotation}{Annotation}] [\href{https://github.com/nmrenyi/extraire-tesseract-openai}{Extraction}] [\href{https://github.com/nmrenyi/rosenwald-benchmark}{Benchmark}] [\href{https://github.com/nmrenyi/rosenwald-extraction}{Données}]
        \end{twocolentry}
        \vspace{0.10 cm}
        \begin{onecolentry}
        \begin{highlights}

    \item Collaboration avec des historien·ne·s du Laboratory for the History of Science and Technology (EPFL) et de l'Institut des humanit\'es en m\'edecine (Lausanne).
    \item Conception d'un cadre d'annotation « double triangle » combinant LLMs et annotation humaine pour une annotation scalable et fiable, réduisant de >50\% l'effort d'annotation humaine tout en maintenant une précision >99\%. \textbf{Prépublication : \href{https://arxiv.org/abs/2605.25781}{arXiv:2605.25781}}.
    \item Construction d'un jeu de données de référence de 2 600+ entrées de médecins des Guides Rosenwald à l'aide du cadre proposé, pour évaluer l'extraction d'informations structurées à partir de documents historiques.
    \item Comparaison de quatre paradigmes d'extraction sur le benchmark et application de la meilleure approche (image + texte OCR) pour extraire 577 000+ entrées de médecins et d'officiels de santé des Guides Rosenwald (1887--1922).
    \item Découverte de 3 700+ entrées de femmes médecins (1887--1922), permettant de futures études sur la pratique médicale féminine.
\end{highlights}
        \end{onecolentry}
        \vspace{0.2 cm}

        \begin{twocolentry}{
            fév. 2026 - juin 2026
        }
            \textbf{Agent conversationnel médical embarqué pour les sages-femmes, \href{https://www.d-tree.org/}{\textit{D-tree}}, Zanzibar, Tanzanie}

            [\href{https://github.com/nmrenyi/mamai}{Code}] [\href{https://arxiv.org/abs/2606.29580}{Prépublication système}] [\href{https://arxiv.org/abs/2606.29467}{Prépublication benchmarks}] [\href{https://nmrenyi-mamai-demo.hf.space}{Démo en ligne}] [\href{https://youtu.be/M_Kruluel28}{Vidéo de démo}]
        \end{twocolentry}
        \vspace{0.10 cm}
        \begin{onecolentry}
            \begin{highlights}
                \item Développement d'un assistant médical RAG entièrement embarqué et hors ligne---un système de recherche EmbeddingGemma-300M et un générateur Google Gemma 4 E4B (int4) tournant sur Android grand public---répondant aux questions cliniques des sages-femmes avec des réponses sourcées, sans connexion et en conservant toutes les données patient sur l'appareil
                \item Constitution d'un \href{https://github.com/nmrenyi/mamai-medical-guidelines}{corpus de directives pour le RAG} de 87 documents de référence (OMS, NICE, MSF et directives nationales de Zanzibar/Tanzanie), découpé en 63 650 passages embarqués
                \item Conception et conduite d'une évaluation par couches du système, couvrant la qualité des réponses de bout en bout, le système de recherche et le générateur mesurés isolément, et la latence embarquée selon la longueur d'entrée et le backend CPU/GPU
                \item Développement d'un \href{https://github.com/nmrenyi/mamai-eval}{harnais d'évaluation} ouvert sur deux benchmarks (\href{https://huggingface.co/datasets/nmrenyi/mamabench}{mamabench} et \href{https://huggingface.co/datasets/nmrenyi/mamaretrieval}{mamaretrieval}), notés par un juge LLM validé contre les grilles de médecins de HealthBench
            \end{highlights}
        \end{onecolentry}


        % \begin{twocolentry}{
        %     May 2021 - Oct 2021
        % }
        %     \textbf{Development and Maintenance of Recommender System in Legal Case Retrieval Platform, Tsinghua University}\end{twocolentry}

        % \vspace{0.10 cm}
        % \begin{onecolentry}
        %     \begin{highlights}
        %         \item Maintained the recommendation module in the Legal Case Retrieval Platform in WeChat Mini-Program
        %         \item Tracked and analyzed the statistics of the recommender system for user behavior analysis
        %     \end{highlights}
        % \end{onecolentry}
        % \vspace{0.20 cm}
